\documentclass[conference]{IEEEtran}

\usepackage{amsmath, amssymb, amsthm}

% Define theorem styles for amsthm to avoid compilation errors
\newtheorem{theorem}{Theorem}
\newtheorem{definition}{Definition}

\title{TLBSS: A TITOS-Integrated Tri-Lateral Binary System Substrate for Coherence-Driven State Dynamics and Invariant-Constrained Measurement}

\author{
    \IEEEauthorblockN{Anonymous Framework}
    \IEEEauthorblockA{Systems Theory Division}
    \IEEEauthorblockA{Obienova}
}

\begin{document}

\maketitle

\begin{abstract}
We present the Tri-Lateral Binary System Substrate (TLBSS), a generative framework defined over inertial and entropic fields producing nonlinear coherence structures. We introduce TITOS (Tri-Lateral Invariant Test Operator System) as a reflexive measurement and closure layer governing invariant-constrained observability. The framework defines coherence emergence, invariant stability, resonance response, collapse modes, and empirical falsifiability conditions within a unified formal system.
\end{abstract}

\section{Introduction}
TLBSS defines a generative state space over two primitive fields: inertia and entropy. TITOS provides a constraint system ensuring that only invariant-preserving operations are admissible, forming a closed epistemic loop between generation, measurement, and validation.

\section{Primitive Structure}

\subsection{Axioms}

\begin{equation}
M \in \mathbb{R}^+, \quad E \in \mathbb{R}^+
\end{equation}

\begin{equation}
\mathcal{S} = \mathcal{S}(M,E)
\end{equation}

\subsection{Coherence Operator}

\begin{equation}
C = M^2 E
\end{equation}

\section{Invariant Structure}

\begin{definition}
The invariant set is defined as:
\begin{equation}
I_C = \mathrm{Inv}(C)
\end{equation}
\end{definition}

\section{TITOS Measurement Layer}

\subsection{Test Operators}

\begin{equation}
T \subseteq I_C
\end{equation}

\subsection{Measurement Error Functional}

\begin{equation}
\mathcal{E}_T = d(T, I_C)
\end{equation}

\subsection{Resonance Operator}

\begin{equation}
R = \mathcal{R}(C, T)
\end{equation}

\begin{equation}
R^* = R(1 - \mathcal{E}_T)
\end{equation}

\subsection{Meta-Invariant Stability}

\begin{equation}
\mathcal{M}_I = \exp(-\|I_C(t) - I_C(t-1)\|)
\end{equation}

\subsection{Global Validity Functional}

\begin{equation}
V = \mathcal{E}_T + (1 - \mathcal{M}_I) + \Delta C
\end{equation}

\section{Closure Theorem}

\begin{theorem}
A TLBSS system is stable if and only if:
\begin{equation}
V < \kappa
\end{equation}
\end{theorem}

\begin{proof}
If $\mathcal{E}_T \rightarrow 0$, $T \subseteq I_C$ holds.  \\
If $\mathcal{M}_I \rightarrow 1$, invariants remain stable.  \\
If $\Delta C \rightarrow 0$, coherence is stationary.  \\
Thus $V \rightarrow 0$, satisfying the stability condition.
\end{proof}

\section{Collapse Modes}

\begin{itemize}
\item Mode I: $C \rightarrow 0$
\item Mode II: $\frac{dI_C}{dt} \neq 0$
\item Mode III: $T \notin I_C$
\item Mode IV: $\mathcal{M}_I \rightarrow 0$
\end{itemize}

\section{Phase Structure}

\begin{itemize}
\item Void Regime
\item Turbulence Regime
\item Rigidity Regime
\item Coherence Zone
\item Resonance Band
\item Meta-Drift Zone
\item Collapse Regime
\end{itemize}

\section{Empirical Mapping}

\begin{equation}
(M,E) \rightarrow \text{observable system variables}
\end{equation}

\begin{equation}
C_{\text{obs}} = M_{\text{obs}}^2 E_{\text{obs}}
\end{equation}

\section{Falsifiability Conditions}

TLBSS is falsified if:

\begin{itemize}
\item $C_{\text{obs}}$ is non-stationary under identical conditions
\item $I_{\text{obs}}$ is non-reproducible
\item $T_{\text{obs}} \notin I_{\text{obs}}$
\item $\mathcal{M}_I \rightarrow 0$ without external forcing
\end{itemize}

\section{Conclusion}

TLBSS defines a tri-field generative substrate governed by nonlinear coherence formation. TITOS imposes invariant-constrained measurement and closure conditions, producing a self-consistent generative--measurement--falsification system.

\end{document}
